\documentclass[a4paper,12pt,oneside]{article}
\makeatletter
\def\input@path{{../../}}
\makeatother
\usepackage{discrete-practice}

\begin{document}

\PracticeHeader{2}{}{ДКА, НКА, регулярные выражения}

\Task{Задание 1}
Пусть $L$ -- формальный язык. Докажите, что $(L^*)^*=L^*$.

\Task{Задание 2}
Пусть $R$ и $S$ -- языки. Докажите или опровергните, что
\[
  (R\mid S)^* = R^*\mid S^*,
  \qquad (R\mid S)^*=(R^*S^*)^*.
\]

\Task{Задание 3}
Рассмотрим ДКА со следующей таблицей переходов.
Опишите неформально язык, допустимый данным ДКА.
\begin{center}
  \includegraphics[width=.49\linewidth]{Practice_automata_1.png}\hfill
  \includegraphics[width=.47\linewidth]{Practice_automata_2.png}
\end{center}

\Task{Задание 4}
Преобразуйте следующий НКА в эквивалентный ДКА.
\begin{center}
  \includegraphics[width=.56\linewidth]{Practice_automata_3.png}
\end{center}

\clearpage

\Task{Задание 5}
Постройте регулярные выражения для следующих языков над алфавитом
$\Sigma=\{0,1\}$.
\begin{alphlist}
  \item[\textbf{а)}] $L_1=\{w\in\Sigma^*\mid$ после каждого 0 в слове $w$ следует $1\}$.
  \item[\textbf{б)}] $L_2=\{w\in\Sigma^*\mid$ вторая или третья буква с конца $w$ это $1\}$.
  \item[\textbf{в)}] $L_3=\{w\in\Sigma^*\mid w$ состоит из чередующихся 0 и 1$\}$.
  \item[\textbf{г)}] $L_4=\{w\in\Sigma^*\mid |w|\ge 2$ и $w$ содержит чётное число нулей$\}$.
\end{alphlist}

\Task{Задание 6}
Постройте регулярное выражение, которое будет распознавать число, которое
делится на четыре. Алфавит: $\Sigma=\{0,1,2,3,4,5,6,7,8,9\}$.

\Task{Задание 7}
Рассмотрите следующий $\varepsilon$-НКА:
\begin{center}
  \includegraphics[width=.7\linewidth]{Practice_automata_4.png}
\end{center}
\begin{alphlist}
  \item[\textbf{а)}] Найдите $\varepsilon$-замыкание каждого из состояний.
\end{alphlist}
{\small\itshape В исходном листке условие пункта б) обрывается.}

\end{document}
